\documentclass[12pt]{article}
\usepackage{amsmath,amssymb}
\usepackage{physics}
\usepackage{empheq}
\usepackage[most]{tcolorbox}
\usepackage{graphicx}
\usepackage{xcolor}
\usepackage{enumitem}
\usepackage{tikz}
\usetikzlibrary{arrows.meta,calc,patterns,decorations.pathmorphing}
\providecommand{\br}[1]{\left(#1\right)}
\providecommand{\sbr}[1]{\left[#1\right]}
\providecommand{\cbr}[1]{\left\{#1\right\}}
\providecommand{\mc}[1]{\mathcal{#1}}
\providecommand{\bk}[2]{\left\langle#1\middle|#2\right\rangle}
\providecommand{\kb}[2]{\left|#1\middle\rangle\middle\langle#2\right|}
\providecommand{\com}[2]{\left[#1,#2\right]}
\providecommand{\mat}[1]{\begin{bmatrix}#1\end{bmatrix}}
\providecommand{\R}{\mathbb{R}}
\providecommand{\C}{\mathbb{C}}
\newtcbox{\mathbox}{nobeforeafter,colback=white,colframe=black,boxrule=0.5pt,arc=3pt}
% ==== Panopticon · Format (change it in the Format panel, or edit here) ====
% panopticon-format: {"paper": "a4", "margin_top": 20, "margin_bottom": 20, "margin_left": 20, "margin_right": 20, "size": "12pt", "font": "cm", "spacing": 1.15, "paragraphs": "space", "head_l": "\\textit{ECE 105 Tutorial 3}", "head_c": "", "head_r": "", "foot_l": "", "foot_c": "{page}", "foot_r": "", "head_rule": true, "theorems": "boxes", "colors": {}, "air": "comfortable"}
\usepackage[a4paper,top=20mm,bottom=20mm,left=20mm,right=20mm,headheight=15pt]{geometry}
\usepackage{setspace}\setstretch{1.15}
\setlength{\parskip}{0.7em plus 0.2em minus 0.1em}\setlength{\parindent}{0pt}
\setlength{\jot}{0.6em}\allowdisplaybreaks[2]
\makeatletter\g@addto@macro\normalsize{\setlength\abovedisplayskip{0.9em plus 0.3em minus 0.2em}\setlength\belowdisplayskip{0.9em plus 0.3em minus 0.2em}\setlength\abovedisplayshortskip{0.4em plus 0.2em}\setlength\belowdisplayshortskip{0.7em plus 0.2em minus 0.1em}}\makeatother
\IfFileExists{microtype.sty}{\usepackage{microtype}}{}
\usepackage{fancyhdr}\pagestyle{fancy}\fancyhf{}
\lhead{\textit{ECE 105 Tutorial 3}}
\cfoot{\thepage}
\renewcommand{\headrulewidth}{0.4pt}
\fancypagestyle{plain}{\fancyhf{}\cfoot{\thepage}\renewcommand{\headrulewidth}{0pt}}
\usepackage{panopticon-thm}  % boxed theorems, like the viewer
% ==== end Format ====
\makeatletter\providecommand{\websvg}{}\renewcommand{\websvg}[1]{\filename@parse{#1}\IfFileExists{\filename@area\filename@base.pdf}{\begingroup\centering\includegraphics[width=0.75\linewidth]{\filename@area\filename@base.pdf}\par\endgroup}{\fbox{\small figure on the website: \texttt{\detokenize{#1}}}}}\makeatother  % Panopticon: \websvg prints the .pdf beside each .svg
\usepackage[hidelinks]{hyperref}

\title{ECE 105 Tutorial 3}
\author{Ritik Kotak}
\date{}
% your extra \newcommand lines go here (one per line)
\newcounter{problem}
\newenvironment{problem}[1][]{\refstepcounter{problem}\par\medskip\noindent\textbf{Problem \theproblem#1.}\ }{\par}
\newenvironment{solution}{\par\noindent\textit{Solution.}\ }{\hfill\rule{0.55em}{0.55em}\par}

\begin{document}
\maketitle

In this tutorial we apply Newton's laws first to systems of connected bodies and then to circular motion. Every problem below rests on the three laws, so we begin by restating them in the form in which we will use them.

\begin{textbox}[Newton's laws, the working version]
\begin{enumerate}[label=\arabic*.]
  \item A body on which the net force is zero moves with constant velocity (in an inertial frame).
  \item The net force on a body equals its mass times its acceleration, $\sum \vec F = m\vec a$. This law is applied to \emph{one body at a time}, and one component at a time.
  \item If body $A$ exerts a force on body $B$, then $B$ exerts a force on $A$ of equal magnitude and opposite direction. These two forces act on \emph{different} bodies, so they never appear on the same free-body diagram and can never cancel each other.
\end{enumerate}
We will also need two kinds of friction. When two surfaces slide over each other, kinetic friction opposes the sliding and has a fixed magnitude, $f_k = \mu_k N$. When they do not slide, static friction takes \emph{whatever value is needed} to prevent sliding, up to a maximum: $f_s \leq \mu_s N$.
\end{textbox}

Problems of this kind can be solved reliably with the following procedure.
\begin{enumerate}[label=(\roman*)]
  \item Isolate each body and draw its free-body diagram (FBD), showing only the forces that act \emph{on} that body.
  \item Choose axes along and perpendicular to the direction of each body's acceleration. The choice may differ from one body to another.
  \item Write Newton's second law for each body along each axis.
  \item Relate the motions of the bodies through the constraints. For example, bodies joined by a taut, inextensible string have accelerations of equal magnitude.
  \item Solve the resulting equations.
\end{enumerate}

\newpage
\section*{Part I: Newton's laws}
%=====================================================================
\begin{problem}\label{prob:ramp}
Consider the system shown in the figure. Blocks of mass $m_A$ and $m_B$ lie on a rough horizontal table and are connected by a light string. A second light string connects $m_B$ to a block of mass $m_C$, which rests on a rough incline of angle $\varphi$. This string passes over a massless, frictionless pulley at the edge of the table and runs parallel to the incline. Both strings are inextensible. For every block, the coefficients of static and kinetic friction with the surface are $\mu_s$ and $\mu_k$ respectively. Neglect air resistance.

\begin{center}
\websvg{figures/tikz-fef0f87a.svg}
\end{center}

\begin{enumerate}[label=(\alph*)]
  \item The string connecting $m_B$ to $m_C$ exerts a force of magnitude $T_2$ on $m_B$, directed towards the pulley. What force does $m_B$ exert on this string? Give its magnitude and direction.
  \item The system is initially at rest. The string connecting $m_B$ to $m_C$ is then cut. Describe what happens to each block immediately after the string is cut. Must $m_C$ slide down the incline? Justify your answers using the forces acting on each block.
\end{enumerate}
For parts (c) to (e), assume the strings are intact and the system is sliding, with $m_C$ moving down the incline and $m_A$ and $m_B$ moving towards the pulley.
\begin{enumerate}[label=(\alph*),resume]
  \item Draw a complete free-body diagram for each of the three blocks.
  \item Using your free-body diagrams, write Newton's second law along the direction of motion for each block. Hence find the acceleration $a$ of the system in terms of $m_A$, $m_B$, $m_C$, $\mu_k$, $\varphi$ and $g$.
  \item Suppose the system slides with constant velocity. Find an equation that determines $\varphi$.
\end{enumerate}
\end{problem}

\begin{solution}

\textbf{(a)} By Newton's third law, if the string exerts a force on $m_B$, then $m_B$ exerts a force on the string of equal magnitude and opposite direction. The block therefore pulls on the string with a force of magnitude $T_2$, directed horizontally away from the pulley. This force acts on the string, not on the block, so it does not appear in the free-body diagram of $m_B$. (What is the third law partner of the weight of $m_C$?)

\textbf{(b)} Once the string is cut, it no longer has anything to pull against, and its tension immediately drops to zero.

Consider first $m_A$ and $m_B$. With $T_2 = 0$, nothing pulls $m_B$ towards the pulley, so the string joining the two blocks also goes slack and $T_1 = 0$. Static friction acts only to oppose a \emph{tendency} to slide. Since no horizontal force now tends to move either block, the friction on each drops to zero, and both blocks remain at rest.

Now consider $m_C$. Before the cut, the string pulled it \emph{up} the incline and helped to hold it in place. After the cut, static friction alone must hold it. The component of gravity along the incline is $m_C g\sin\varphi$, while the largest friction force the surface can supply is $\mu_s N_C = \mu_s m_C g\cos\varphi$. The block therefore begins to slide only if
\begin{flalign}
\tan\varphi &> \mu_s &&
\end{flalign}
Hence $m_C$ need not slide. Whether it does depends only on the angle of the incline and the coefficient of static friction, not on its mass. Notice that removing a force (here, the tension) can set a body in motion.

\textbf{(c)} We take the positive direction to be towards the pulley for $m_A$ and $m_B$, and down the incline for $m_C$. Since all three blocks are sliding, the friction on each is kinetic and opposes its motion.

\begin{center}
\websvg{figures/tikz-d0440d65.svg}
\end{center}

Three features of these diagrams deserve attention. First, $m_B$ is acted on by \emph{two} tensions: $T_2$ pulls it towards the pulley and $T_1$ pulls it back towards $m_A$. Second, the tension $T_2$ on $m_C$ acts \emph{up} the incline, because the string runs parallel to the incline; a massless, frictionless pulley changes the direction of the string but not its tension. Third, no ``$ma$'' arrow appears on any diagram, since a free-body diagram shows only forces exerted by other bodies.

\textbf{(d)} None of the blocks accelerates perpendicular to its surface, so the normal forces are $N_A = m_A g$, $N_B = m_B g$ and $N_C = m_C g\cos\varphi$. Because the strings are inextensible, all three blocks have the same magnitude of acceleration $a$ (Why?). Newton's second law along the direction of motion of each block gives
\begin{flalign}
m_A:&\quad T_1 - \mu_k m_A g = m_A a &&\\
m_B:&\quad T_2 - T_1 - \mu_k m_B g = m_B a &&\\
m_C:&\quad m_C g\sin\varphi - T_2 - \mu_k m_C g\cos\varphi = m_C a &&
\end{flalign}
When the three equations are added, the tensions cancel in pairs, since they are internal forces of the system. Solving for $a$ gives
\begin{empheq}[box=\mathbox]{align}
a = \frac{m_C\br{\sin\varphi - \mu_k\cos\varphi} - \mu_k\br{m_A + m_B}}{m_A + m_B + m_C}\; g
\end{empheq}
The numerator is the net driving force (the component of $m_C$'s weight along the incline minus all of the friction forces), and the denominator is the total mass. In other words, the result is simply $\sum F_{\text{ext}} = Ma$ applied to the system as a whole.

\textbf{(e)} If the system slides with constant velocity, its acceleration is zero. Setting $a = 0$ gives
\begin{empheq}[box=\mathbox]{align}
m_C\br{\sin\varphi - \mu_k\cos\varphi} = \mu_k\br{m_A + m_B}
\end{empheq}
As a check, setting $m_A = m_B = 0$ gives $\tan\varphi = \mu_k$, the familiar condition for a single block to slide down an incline at constant speed.
\end{solution}

\newpage
%=====================================================================
\begin{problem}\label{prob:wedge}
A block of mass $m$ rests on the smooth inclined face of a wedge of mass $M$ and angle $\theta$. The wedge stands on a smooth horizontal floor. A horizontal force $F$ is applied to the vertical face of the wedge, and its magnitude is chosen so that the block does not slide relative to the wedge.

\begin{center}
\websvg{figures/tikz-7a72a65c.svg}
\end{center}

\begin{enumerate}[label=(\alph*)]
  \item Draw complete free-body diagrams of the block and of the wedge. Identify every third law pair that connects the two diagrams.
  \item In which direction does the block accelerate? Use this to choose suitable axes, and find the acceleration $a$ and the normal force $N$ on the block.
  \item Find the magnitude of $F$.
  \item For a block on a fixed incline, $N = mg\cos\theta$. Explain why $N$ is different here. Is it larger or smaller?
  \item If $F$ is made slightly larger than the value found in (c), in which direction does the block slide relative to the wedge?
\end{enumerate}
\end{problem}

\begin{solution}

\textbf{(a)} Since the incline is smooth, only two forces act on the block: its weight $mg$ and the normal force $N$ exerted by the incline, which is perpendicular to the incline. Four forces act on the wedge: the applied force $F$, its weight $Mg$, the normal force $N_{\text{floor}}$ from the floor, and the force exerted by the block. By Newton's third law, this last force has the same magnitude $N$ as the normal force on the block, and it is directed perpendicular to the incline, \emph{into} the wedge.

\begin{center}
\websvg{figures/tikz-3bc72818.svg}
\end{center}

The two forces labelled $N$ form the only third law pair connecting the two diagrams. Note that the weight of the block and the normal force on it are \emph{not} a third law pair, since they act on the same body. (What is the third law partner of $N_{\text{floor}}$?)

\textbf{(b)} Since the block does not slide relative to the wedge, it shares the wedge's acceleration, which is \emph{horizontal} and in the direction of $F$. It is therefore convenient to use horizontal and vertical axes rather than axes along and perpendicular to the incline, because with this choice only one of the two equations contains the acceleration. Since the block has no vertical acceleration,
\begin{flalign}
\text{vertical:}&\quad N\cos\theta - mg = 0 &&\\
\text{horizontal:}&\quad N\sin\theta = ma &&
\end{flalign}
\begin{empheq}[box=\mathbox]{align}
a = g\tan\theta, \qquad N = \frac{mg}{\cos\theta}
\end{empheq}

\textbf{(c)} Since the block and the wedge move together, we may treat them as a single body of mass $M + m$. The forces between them are internal and do not enter, so
\begin{flalign}
F &= (M + m)\, a = (M + m)\, g\tan\theta &&
\end{flalign}
We can confirm this by applying Newton's second law to the wedge alone, $F - N\sin\theta = Ma$, which gives the same result.

\textbf{(d)} The normal force $N = mg/\cos\theta$ is \emph{larger} than $mg\cos\theta$. The familiar expression $N = mg\cos\theta$ is not a general law; it holds only when the block has no acceleration perpendicular to the incline. Here the block accelerates horizontally, and this acceleration has a component $a\sin\theta$ directed away from the surface of the incline. Newton's second law perpendicular to the incline then gives
\begin{flalign}
N - mg\cos\theta &= m\, a\sin\theta \quad\Longrightarrow\quad N = mg\cos\theta + mg\frac{\sin^2\theta}{\cos\theta} = \frac{mg}{\cos\theta} &&
\end{flalign}
The normal force always takes whatever value the motion requires, so it should be found from Newton's second law rather than from a memorized formula.

\textbf{(e)} A larger $F$ requires a larger horizontal acceleration of the block. Only the horizontal component $N\sin\theta$ can supply this, so $N$ must increase. Its vertical component $N\cos\theta$ then exceeds $mg$, and the block accelerates upwards relative to the wedge; that is, it slides \emph{up} the incline. Conversely, if $F$ is made smaller, the block slides down the incline.
\end{solution}

%=====================================================================
\begin{problem}[$^\ast$]\label{prob:freewedge}
Consider again the wedge of Problem~\ref{prob:wedge}, of mass $M$ and angle $\theta$, standing on a smooth horizontal floor, but now with no applied force. A block of mass $m$ is released from rest on its smooth inclined face. As the block slides down the incline, the wedge slides along the floor.

\begin{center}
\websvg{figures/tikz-607ad700.svg}
\end{center}

\begin{enumerate}[label=(\alph*)]
  \item Draw free-body diagrams of the block and of the wedge. In which direction does the wedge accelerate, and why?
  \item Let $A$ be the magnitude of the wedge's acceleration, and let $b$ be the magnitude of the block's acceleration \emph{relative to the wedge}. Explain why the block's acceleration relative to the wedge must be directed along the incline, and write the block's acceleration in the ground frame in terms of $A$ and $b$.
  \item Find $A$ and the normal force $N$ between the block and the wedge. Is $N$ larger or smaller than $mg\cos\theta$? Explain physically.
  \item Show that, in the ground frame, the block moves along a straight line, and find the angle this line makes with the horizontal. Is the path steeper or shallower than the incline?
  \item Find the normal force exerted by the floor on the wedge. Why is it less than $(M + m)g$?
  \item Check your results in the limits $M \to \infty$ and $\theta \to 90^\circ$.
\end{enumerate}
\end{problem}

\begin{solution}

\textbf{(a)} Two forces act on the block: its weight $mg$ and the normal force $N$, perpendicular to the incline. Three forces act on the wedge: its weight $Mg$, the normal force $N_{\text{floor}}$ from the floor (vertical, since the floor is smooth), and the force exerted by the block, which by Newton's third law has magnitude $N$ and is directed perpendicular to the incline, \emph{into} the wedge.

\begin{center}
\websvg{figures/tikz-ef38dfd0.svg}
\end{center}

The only horizontal force on the wedge is the horizontal component $N\sin\theta$ of the block's push. This component points away from the side towards which the block slides, so the wedge accelerates \emph{backwards}, opposite to the horizontal motion of the block.

\textbf{(b)} The block remains in contact with the incline: it neither leaves the surface nor penetrates it. Relative to the wedge, it can therefore move only along the surface, and its relative acceleration $b$ is directed down the incline. Its acceleration in the ground frame is the sum of the wedge's acceleration and this relative acceleration, as shown in the right-hand diagram. Taking the $x$-axis in the direction in which the block slides and the $y$-axis vertically upwards,
\begin{flalign}
\vec a &= \underbrace{(-A,\ 0)}_{\text{wedge}} + \underbrace{(b\cos\theta,\ -b\sin\theta)}_{\text{relative}} &&
\end{flalign}
This step is what makes the problem difficult. The block's acceleration in the ground frame is \emph{not} directed along the incline, so Newton's second law cannot simply be written ``along the incline'' for the block.

\textbf{(c)} There are three unknowns, $A$, $b$ and $N$, and three equations: the horizontal and vertical components of Newton's second law for the block, and the horizontal component for the wedge.
\begin{flalign}
\text{block, } x:&\quad N\sin\theta = m(b\cos\theta - A) &&\\
\text{block, } y:&\quad N\cos\theta - mg = -m b\sin\theta &&\\
\text{wedge, } x:&\quad N\sin\theta = MA &&
\end{flalign}
From the first and third equations, $b = (M+m)A/(m\cos\theta)$. Substituting this, together with $N = MA/\sin\theta$, into the second equation and simplifying gives
\begin{empheq}[box=\mathbox]{align}
A = \frac{mg\sin\theta\cos\theta}{M + m\sin^2\theta}, \qquad N = \frac{Mmg\cos\theta}{M + m\sin^2\theta}
\end{empheq}
and the relative acceleration is $b = (M+m)g\sin\theta/(M + m\sin^2\theta)$.

Since $M + m\sin^2\theta > M$, we find $N < mg\cos\theta$. Physically, the wedge recoils away from the block, so the incline recedes beneath the block and need not push as hard as a fixed incline would. This is the opposite of Problem~\ref{prob:wedge}, in which the wedge was pushed \emph{into} the block and $N$ was larger than $mg\cos\theta$. In both cases, the normal force is determined by the motion rather than by a fixed formula.

\textbf{(d)} Since $A$ and $b$ are both constant, the block's acceleration $\vec a$ is constant. A body that starts from rest with constant acceleration moves in a straight line in the direction of that acceleration. Using part (b) and the relation $b = (M+m)A/(m\cos\theta)$, the components of $\vec a$ are
\begin{flalign}
a_x &= b\cos\theta - A = \frac{M}{m}A, \qquad a_y = -b\sin\theta = -\frac{M+m}{m}A\tan\theta &&
\end{flalign}
The path therefore makes an angle $\beta$ below the horizontal, where
\begin{empheq}[box=\mathbox]{align}
\tan\beta = \frac{M + m}{M}\tan\theta > \tan\theta
\end{empheq}
The path is \emph{steeper} than the incline, because the wedge moves out from under the block. As a check, note that $m a_x = MA$. This must be so, since the horizontal forces on the block and on the wedge form a third law pair and no external horizontal force acts on the system.

\textbf{(e)} The only vertical external forces on the system consisting of the block and the wedge are gravity and the normal force from the floor. The wedge has no vertical acceleration, while the block accelerates downwards at $b\sin\theta$. Newton's second law for the system in the vertical direction gives
\begin{flalign}
(M+m)g - N_{\text{floor}} &= m b\sin\theta \quad\Longrightarrow\quad N_{\text{floor}} = (M+m)g - \frac{m(M+m)g\sin^2\theta}{M + m\sin^2\theta} &&
\end{flalign}
The floor supports less than the total weight because part of the system is accelerating downwards. (Verify that this agrees with $N_{\text{floor}} = Mg + N\cos\theta$, obtained from the free-body diagram of the wedge alone.)

\textbf{(f)} In the limit $M \to \infty$, we find $A \to 0$, $N \to mg\cos\theta$, $b \to g\sin\theta$ and $\beta \to \theta$, which are the results for a fixed incline. In the limit $\theta \to 90^\circ$, we find $A \to 0$, $N \to 0$ and $b \to g$: the block falls freely alongside the vertical face of the wedge without pressing against it.
\end{solution}

\newpage
%=====================================================================
% SOURCE: Irodov, Problems in General Physics, Sec. 1.2 (two touching bars on an incline)
\begin{problem}\label{prob:bars}
Two bars, 1 and 2, are placed in contact on an incline of angle $\alpha$, with bar 1 below bar 2. The masses of the bars are $m_1$ and $m_2$, and their coefficients of friction with the incline are $\mu_1$ and $\mu_2$ respectively, where $\mu_1 > \mu_2$. Assume that the static and kinetic coefficients of friction are equal.

\begin{center}
\websvg{figures/tikz-8193d613.svg}
\end{center}

\begin{enumerate}[label=(\alph*)]
  \item Draw the forces acting along the incline on each bar while the bars slide down together.
  \item Find the common acceleration of the bars and the force that each bar exerts on the other.
  \item Find the minimum angle $\alpha$ at which the bars begin to slide.
  \item What changes if the order of the bars is reversed, so that bar 1 is above bar 2?
\end{enumerate}
\end{problem}

\begin{solution}

\textbf{(a)} We take the direction down the incline as positive. Each bar experiences a component of gravity $mg\sin\alpha$ down the incline and a friction force $\mu mg\cos\alpha$ up the incline, since the normal force on each bar is $mg\cos\alpha$. Bar 2 experiences the smaller friction force per unit mass, so on its own it would accelerate faster than bar 1. Since it is behind bar 1, it pushes on bar 1 with a force $R$ directed down the incline. By Newton's third law, bar 1 pushes back on bar 2 with a force $R$ directed up the incline.

\begin{center}
\websvg{figures/tikz-6052118b.svg}
\end{center}
The bars are drawn slightly apart so that the two contact forces can be seen. In the diagram, $f_1 = \mu_1 m_1 g\cos\alpha$ and $f_2 = \mu_2 m_2 g\cos\alpha$ act up the incline along the contact surfaces, and the pair of forces $R$ (in red) is the third law pair between the bars.

\textbf{(b)} Since the bars remain in contact, they share a common acceleration $a$. Newton's second law for each bar gives
\begin{flalign}
m_1 a &= m_1 g\sin\alpha - \mu_1 m_1 g\cos\alpha + R &&\\
m_2 a &= m_2 g\sin\alpha - \mu_2 m_2 g\cos\alpha - R &&
\end{flalign}
Adding these equations eliminates $R$:
\begin{flalign}
a &= g\sin\alpha - \frac{\mu_1 m_1 + \mu_2 m_2}{m_1 + m_2}\, g\cos\alpha &&
\end{flalign}
The pair therefore moves like a single body whose coefficient of friction is the mass-weighted average $\bar\mu = (\mu_1 m_1 + \mu_2 m_2)/(m_1+m_2)$. Substituting $a$ back into the first equation gives the contact force
\begin{empheq}[box=\mathbox]{align}
R = \frac{m_1 m_2 (\mu_1 - \mu_2)}{m_1 + m_2}\, g\cos\alpha
\end{empheq}
The bars push on each other only because $\mu_1 \neq \mu_2$. If the coefficients were equal, each bar would slide with the same acceleration on its own, and $R$ would be zero.

\textbf{(c)} The bars begin to slide when $a > 0$, which gives the minimum angle
\begin{flalign}
\tan\alpha_{\min} &= \bar\mu = \frac{\mu_1 m_1 + \mu_2 m_2}{m_1 + m_2} &&
\end{flalign}
On its own, bar 2 would begin to slide at $\tan\alpha = \mu_2$, and bar 1 at $\tan\alpha = \mu_1$. Together, they begin to slide at an intermediate angle, because the more slippery bar helps to push the other one down.

\textbf{(d)} If the order is reversed, repeating the calculation of part (b) gives $R < 0$. Since a contact force can only push, this is impossible, and the bars cannot remain in contact. Bar 1, which has the larger coefficient of friction, now lags behind, and the bars separate. Each then slides with its own acceleration $g(\sin\alpha - \mu\cos\alpha)$. In general, a negative contact force in a calculation indicates that the bodies are not actually in contact.
\end{solution}

%=====================================================================
% SOURCE: Irodov, Problems in General Physics, Sec. 1.2 (bar pulled up an incline, minimum tension)
\begin{problem}\label{prob:pull}
A bar of mass $m$ is pulled up a rough incline of angle $\alpha$ by means of a thread. The coefficient of friction between the bar and the incline is $\mu$, and the thread makes an angle $\beta$ with the incline.

\begin{center}
\websvg{figures/tikz-fc0820de.svg}
\end{center}

\begin{enumerate}[label=(\alph*)]
  \item Find the tension in the thread required to pull the bar up the incline at constant velocity, as a function of $\beta$.
  \item For what value of $\beta$ is this tension a minimum? What is the minimum tension?
  \item Explain why pulling at an angle to the incline requires less tension than pulling parallel to it.
\end{enumerate}
\end{problem}

\begin{solution}

\textbf{(a)} Four forces act on the bar: its weight $mg$, the normal force $N$, the tension $T$ at an angle $\beta$ to the incline, and the kinetic friction force $\mu N$ directed down the incline.

\begin{center}
\websvg{figures/tikz-1d97b9e6.svg}
\end{center}

Since the bar moves with constant velocity, the forces balance both perpendicular and parallel to the incline:
\begin{flalign}
\perp:&\quad N + T\sin\beta = mg\cos\alpha &&\\
\parallel:&\quad T\cos\beta = mg\sin\alpha + \mu N &&
\end{flalign}
The first equation shows that the thread lifts part of the bar's weight off the incline, so that $N$ is \emph{less} than $mg\cos\alpha$. Eliminating $N$ between the two equations gives
\begin{flalign}
T(\beta) &= \frac{mg\br{\sin\alpha + \mu\cos\alpha}}{\cos\beta + \mu\sin\beta} &&
\end{flalign}

\textbf{(b)} The tension is a minimum when the denominator is a maximum. Setting $\frac{d}{d\beta}\br{\cos\beta + \mu\sin\beta} = 0$ gives $-\sin\beta + \mu\cos\beta = 0$, and therefore
\begin{empheq}[box=\mathbox]{align}
\tan\beta = \mu, \qquad T_{\min} = \frac{mg\br{\sin\alpha + \mu\cos\alpha}}{\sqrt{1 + \mu^2}}
\end{empheq}
The optimal angle depends only on the coefficient of friction, not on the angle of the incline or the mass of the bar.

\textbf{(c)} Tilting the thread upwards has two competing effects. It reduces the component of the tension along the incline by a factor $\cos\beta$, but it also reduces the normal force and hence the friction. For small $\beta$, the loss is of second order, since $\cos\beta \approx 1 - \beta^2/2$, while the gain is of first order, since $\mu\sin\beta \approx \mu\beta$. A small tilt therefore always helps, and the optimal angle is the one at which the two effects balance.

There is also a geometric way to see why $\tan\beta = \mu$. The normal force $N$ and the friction force $\mu N$ may be combined into a single contact force exerted by the incline. Whatever the value of $N$, this force makes the fixed angle $\arctan \mu$ with the normal. In equilibrium, the weight, the contact force and the tension must form a closed triangle. Since the weight is fixed and the contact force has a fixed direction, the tension is shortest when it is perpendicular to the contact force, which places the thread at an angle $\arctan \mu$ to the incline.
\end{solution}

%=====================================================================
% SOURCE: Irodov, Problems in General Physics, Sec. 1.2 (plank and bar, F = at applied to the bar)
\begin{problem}[$^\ast$]\label{prob:plank}
A plank of mass $m_1$ rests on a smooth horizontal floor, and a bar of mass $m_2$ rests on top of the plank. The coefficient of friction between the bar and the plank is $\mu$. A horizontal force $F = \lambda t$, where $\lambda$ is a constant, is applied to the bar.

\begin{center}
\websvg{figures/tikz-e424740e.svg}
\end{center}

\begin{enumerate}[label=(\alph*)]
  \item At what time $t_0$ does the bar begin to slide relative to the plank?
  \item Find the acceleration $a_1$ of the plank and the acceleration $a_2$ of the bar as functions of time, and sketch both on the same axes.
\end{enumerate}
\end{problem}

\begin{solution}

\textbf{(a)} The free-body diagrams of the bar and the plank are shown below. The two bodies interact through a normal force $N$ and a friction force $f$, and each of these forms a third law pair (shown in red). Drawing the forces on the set-up itself would place both members of each pair on the same contact surface, so the bodies are drawn separately.

\begin{center}
\websvg{figures/tikz-f75eea67.svg}
\end{center}

Since the floor is smooth, the \emph{only} horizontal force on the plank is the friction force $f$ exerted by the bar, and vertically $N = m_2 g$. While the bar and the plank move together, their common acceleration and the friction force on the plank are
\begin{flalign}
a &= \frac{\lambda t}{m_1 + m_2}, \qquad f = m_1 a = \frac{m_1\, \lambda t}{m_1 + m_2} &&
\end{flalign}
As $t$ increases, the static friction force grows so as to keep the plank moving with the bar. However, it cannot exceed $\mu m_2 g$, so the acceleration of the plank can be at most $\mu m_2 g/m_1$. The bar begins to slide when the required friction force reaches this maximum, which happens at
\begin{empheq}[box=\mathbox]{align}
t_0 = \frac{\mu g\, m_2 (m_1 + m_2)}{\lambda\, m_1}
\end{empheq}

\textbf{(b)} For $t \leq t_0$, the two bodies move together with $a_1 = a_2 = \lambda t/(m_1+m_2)$. For $t > t_0$, the friction is kinetic and has the constant magnitude $\mu m_2 g$, so
\begin{flalign}
a_1 &= \frac{\mu m_2 g}{m_1}, \qquad a_2 = \frac{\lambda t - \mu m_2 g}{m_2} &&
\end{flalign}
Once slipping begins, the acceleration of the plank remains constant at its maximum value. The acceleration of the bar increases more rapidly than before, with slope $\lambda/m_2$ instead of $\lambda/(m_1+m_2)$, because the applied force no longer has to accelerate the plank. (Verify that $a_2$ is continuous at $t_0$.)

\begin{center}
\websvg{figures/tikz-b66fa72e.svg}
\end{center}
\end{solution}

%=====================================================================
\begin{problem}\label{prob:monkey}
A monkey of mass $m$ sits in a basket of mass $M$. A light rope passes over a massless, frictionless pulley fixed to the ceiling. One end of the rope is tied to the basket, and the monkey holds the other end.

\begin{center}
\websvg{figures/tikz-fe2d6627.svg}
\end{center}

\begin{enumerate}[label=(\alph*)]
  \item With what force must the monkey pull on the rope to rise at constant speed? With what force must it pull to rise with acceleration $a$?
  \item Find the normal force between the monkey and the floor of the basket. Under what condition is the monkey unable to remain seated while pulling itself up?
  \item Now consider a different arrangement. A monkey of mass $m$ hangs from one end of a rope that passes over the same pulley, and a counterweight of mass $m$ hangs from the other end, so that the system is balanced. The monkey begins to climb the rope. What happens to the counterweight?
\end{enumerate}
\end{problem}

\begin{solution}

\textbf{(a)} We take the monkey and the basket together as our system. Since the rope is massless and the pulley is massless and frictionless, the tension $T$ is the same throughout the rope. The rope pulls upwards on the system at \emph{two} points, once at the basket and once at the monkey's hands, so the total upward force is $2T$.
\begin{center}
\websvg{figures/tikz-a3a06a79.svg}
\end{center}
Newton's second law for the system gives
\begin{flalign}
2T - (m+M)g &= (m+M)a \quad\Longrightarrow\quad T = \frac{(m+M)(g+a)}{2} &&
\end{flalign}
At constant speed ($a = 0$), $T = (m+M)g/2$. By Newton's third law, the monkey pulls on the rope with this same force. The monkey therefore needs to exert only \emph{half} of the combined weight, because the rope supports the system at two points.

\textbf{(b)} Now consider the monkey alone (right-hand diagram). The rope pulls it upwards with force $T$, and the floor of the basket pushes it upwards with force $N$. Newton's second law gives
\begin{flalign}
T + N - mg &= ma \quad\Longrightarrow\quad N = m(g+a) - T = \frac{(m - M)(g+a)}{2} &&
\end{flalign}
Since the floor can only push, we require $N \geq 0$, that is, $m \geq M$. If the basket is heavier than the monkey, the monkey rises off the floor and the basket is left behind. To take the basket with it, the monkey would have to hold on to it, so that the contact force could become a pull.

\textbf{(c)} The tension is the same on both sides of the pulley. Whatever the monkey does, whether it climbs steadily, jerks or stops, it can affect the system only by changing $T$, and $T$ acts in exactly the same way on both bodies:
\begin{flalign}
\text{monkey:}&\quad T - mg = m a_{\text{monkey}} &&\\
\text{weight:}&\quad T - mg = m a_{\text{weight}} &&
\end{flalign}
Hence $a_{\text{monkey}} = a_{\text{weight}}$ at every instant. Since both start from rest, the counterweight rises \emph{exactly} alongside the monkey. If the monkey climbs at speed $u$ relative to the rope, each body rises at speed $u/2$ relative to the ground; the climbing simply draws rope over the pulley.
\end{solution}

%=====================================================================
\begin{problem}\label{prob:pulley}
A block of mass $m_1$ hangs from one end of a light, inextensible string. The string passes over a fixed pulley, down and around a movable pulley, and up to the ceiling, where its other end is attached. A block of mass $m_2$ hangs from the axle of the movable pulley. Both pulleys are massless and frictionless.

\begin{center}
\websvg{figures/tikz-63d04944.svg}
\end{center}

\begin{enumerate}[label=(\alph*)]
  \item Explain why the tension is the same throughout the string. Draw free-body diagrams of $m_1$ and of the movable pulley together with $m_2$.
  \item If $m_1$ moves down a distance $d$, how far does $m_2$ move? Hence find the relation between their accelerations.
  \item Find the accelerations of $m_1$ and $m_2$.
  \item For what ratio $m_2/m_1$ is the system in equilibrium? Explain physically why this ratio is not 1.
  \item Suppose $m_1 = m_2 = m$. In which direction does each block move, and with what acceleration?
\end{enumerate}
\end{problem}

\begin{solution}

\textbf{(a)} Consider any massless object, such as a short segment of string or a massless pulley. If a nonzero net force acted on it, Newton's second law would give it an infinite acceleration. The net force on any massless object must therefore be zero. For a massless string, this means that the forces pulling on its two ends are equal. For a massless, frictionless pulley, it means that the string pulls equally on both sides. Hence a single tension $T$ exists throughout the string.

The movable pulley is supported by \emph{two} segments of string, so the pulley together with $m_2$ experiences an upward force $2T$ and a downward force $m_2 g$. Block $m_1$ experiences an upward force $T$ and a downward force $m_1 g$.

\begin{center}
\websvg{figures/tikz-f5f2bd0a.svg}
\end{center}

\textbf{(b)} We measure positions downwards, letting $y_1$ be the position of $m_1$ and $y_2$ the position of the movable pulley. Apart from the fixed lengths wrapped around the pulleys, the string consists of one segment of length $y_1$ (up to a constant) and two segments of length $y_2$ (up to a constant). Since the total length of the string is fixed,
\begin{flalign}
y_1 + 2y_2 &= \text{const} \quad\Longrightarrow\quad a_1 = -2a_2 &&
\end{flalign}
Thus, if $m_1$ descends a distance $d$, this length of string is shared between the two segments supporting the movable pulley, and $m_2$ rises a distance $d/2$.

\textbf{(c)} Taking the downward direction as positive for both blocks, Newton's second law gives
\begin{flalign}
m_1 g - T &= m_1 a_1 &&\\
m_2 g - 2T &= m_2 a_2 = -\tfrac12 m_2 a_1 &&
\end{flalign}
Multiplying the first equation by two and subtracting the second eliminates $T$, and we obtain
\begin{empheq}[box=\mathbox]{align}
a_1 = \frac{2(2m_1 - m_2)}{4m_1 + m_2}\, g, \qquad a_2 = -\frac{2m_1 - m_2}{4m_1 + m_2}\, g
\end{empheq}

\textbf{(d)} The system is in equilibrium when $a_1 = 0$, that is, when $m_2 = 2m_1$. The movable pulley halves the force required: the weight of $m_1$ determines $T$, and the two segments of string together support $m_2$ with a force $2T$. The cost is distance. By part (b), $m_2$ moves only half as far as $m_1$, so the work $m_1 g d$ done by gravity on $m_1$ equals the work $m_2 g (d/2)$ done against gravity on $m_2$. This trade-off between force and distance underlies every pulley system.

\textbf{(e)} With $m_1 = m_2 = m$, we find $a_1 = \tfrac{2}{5}g$ downwards and $a_2 = \tfrac15 g$ upwards. Thus $m_1$ descends even though the two masses are equal, because, acting through the movable pulley, $m_2$ pulls on the string like a mass of only $m_2/2$. Note also that the denominator is $4m_1 + m_2$ rather than $m_1 + m_2$. Since $m_2$ moves at half the speed of $m_1$, its contribution to the inertia of the system is only one quarter as large. (Why one quarter?)
\end{solution}

\newpage
\section*{Part II: Circular motion}
%=====================================================================
\begin{textbox}[Uniform circular motion]
A body moving in a circle of radius $r$ at constant angular velocity $\omega$, and hence at constant speed $v = \omega r$, is nevertheless accelerating, because the direction of its velocity changes continuously. Its acceleration is directed towards the centre of \emph{its} circle and has magnitude
\begin{flalign}
a_c = \frac{v^2}{r} = \omega^2 r &&
\end{flalign}
The term ``centripetal force'' does not refer to a new kind of force, and it never appears as a separate arrow on a free-body diagram. It is simply the name for the component of the net force directed towards the centre, which may be provided by tension, a normal force, gravity, friction, or a combination of these. Newton's second law in the radial direction reads
\begin{flalign}
\sum F_{\text{towards centre}} = m\omega^2 r &&
\end{flalign}
The centre in question is the centre of the circle that the body actually traces, which is not always the most obvious centre in the diagram.
\end{textbox}

\begin{problem}\label{prob:cone}
A small ball of mass $m$ is attached to a light string of length $L$, whose upper end is fixed to the top of a vertical shaft. The shaft rotates with constant angular velocity $\omega$, and the ball moves with it in a horizontal circle, with the string making an angle $\theta$ with the vertical.

\begin{center}
\websvg{figures/tikz-bf3c36e5.svg}
\end{center}

\begin{enumerate}[label=(\alph*)]
  \item Draw the free-body diagram of the ball. Where is the centre of the circle on which the ball moves, and what is its radius $r$? In which direction does the net force on the ball point?
  \item Write Newton's second law for the ball in the vertical direction and in the horizontal (radial) direction.
  \item Find the tension $T$ and the angle $\theta$ in terms of $m$, $g$, $L$ and $\omega$.
  \item Show that the string can make a nonzero angle with the shaft only if $\omega$ exceeds a critical value $\omega_c$, and find $\omega_c$. What does the ball do when $\omega < \omega_c$?
  \item Take $L = 0.50$ m and $\omega = 2\omega_c$. Find $\theta$, and express $T$ as a multiple of $mg$. Can $\theta$ ever reach $90^\circ$?
\end{enumerate}
\end{problem}

\begin{solution}

\textbf{(a)} Only two bodies interact with the ball: the string and the Earth. Accordingly, two forces act on it: the tension $T$, directed along the string towards the pivot, and its weight $mg$, directed vertically downwards.

\begin{center}
\websvg{figures/tikz-5026a35b.svg}
\end{center}

The ball moves in a \emph{horizontal} circle, so the centre of this circle lies on the shaft at the same height as the ball, not at the pivot. The radius of the circle is
\begin{flalign}
r = L\sin\theta &&
\end{flalign}
The ball's acceleration has magnitude $\omega^2 r$ and is directed horizontally towards the shaft. The net force on the ball must therefore also be horizontal and directed towards the shaft; in particular, it does \emph{not} point along the string. (Why can the net force have no vertical component?)

\textbf{(b)} The ball has no vertical acceleration, and its horizontal acceleration is $\omega^2 r$ towards the shaft. Newton's second law therefore gives
\begin{flalign}
T\cos\theta - mg &= 0 && \label{eq:vert}\\
T\sin\theta &= m\omega^2 L\sin\theta && \label{eq:horiz}
\end{flalign}
The vertical component of the tension supports the weight of the ball, while the horizontal component provides the centripetal force.

\textbf{(c)} Equation \eqref{eq:horiz} contains a factor of $\sin\theta$ on both sides. Dividing by this factor, on the assumption that $\sin\theta \neq 0$, gives
\begin{flalign}
T = m\omega^2 L &&
\end{flalign}
and substituting this into \eqref{eq:vert} gives
\begin{empheq}[box=\mathbox]{align}
\cos\theta = \frac{g}{\omega^2 L}
\end{empheq}
Some care is needed here, because dividing by $\sin\theta$ discarded a solution. The value $\theta = 0$ also satisfies \eqref{eq:horiz}: the ball hangs vertically and rotates in place on the axis, with $T = mg$. Whenever we divide by a quantity, we should ask what happens when that quantity is zero.

\textbf{(d)} Since $\cos\theta \leq 1$, a solution with $\theta > 0$ exists only if
\begin{flalign}
\frac{g}{\omega^2 L} \leq 1 \quad\Longrightarrow\quad \omega \geq \omega_c = \sqrt{\frac{g}{L}} &&
\end{flalign}
For $\omega < \omega_c$, the only solution is $\theta = 0$, and the ball simply hangs against the shaft and rotates with it. Physically, at low angular velocity the ball requires only a small inward force, but any tilt of the string would produce an inward force that is too large. A tilted string becomes consistent with the motion only once the rotation is fast enough.

Notice also that the mass $m$ cancels throughout. A heavier ball requires a larger centripetal force, but it also increases the tension in the string by exactly the same factor.

\textbf{(e)} The critical angular velocity is $\omega_c = \sqrt{9.8/0.50} \approx 4.4$ rad/s, so $\omega \approx 8.9$ rad/s. Then
\begin{flalign}
\cos\theta = \frac{g}{4\omega_c^2 L} = \frac14 \quad\Longrightarrow\quad \theta \approx 75.5^\circ, \qquad T = m(2\omega_c)^2 L = 4mg &&
\end{flalign}
The angle $\theta$ can never reach $90^\circ$, however fast the shaft rotates. At $90^\circ$ the string would be horizontal, with no vertical component to support the weight, and \eqref{eq:vert} would require an infinite tension.
\end{solution}

\newpage
%=====================================================================
\begin{problem}[$^\ast$]\label{prob:bead}
A bead of mass 10 mg can slide freely on a frictionless circular wire loop of diameter 20 cm. The loop rotates about its vertical diameter with angular velocity $\omega$. If $\omega$ is less than a certain critical value $\omega_c$, the bead remains at the bottom of the rotating loop. If $\omega > \omega_c$, the bead moves out to an angle $\theta$, measured from the downward vertical through the centre of the loop.

\begin{center}
\websvg{figures/tikz-e2c7f1c9.svg}
\end{center}

\begin{enumerate}[label=(\alph*)]
  \item Draw the free-body diagram of the bead when it is at angle $\theta$. Where is the centre of the circle on which the bead moves, and what is its radius?
  \item Find $\theta$ as a function of $\omega$, and hence find $\omega_c$ in revolutions per minute.
  \item At what angular velocity, in revolutions per minute, is the bead at $\theta = 30^\circ$?
  \item Do your answers depend on the mass of the bead?
  \item[(e)$^\ast$] By considering the forces along the wire, explain why the bead cannot remain at any angle $\theta > 0$ when $\omega < \omega_c$, and why it does not remain at the bottom when $\omega > \omega_c$.
\end{enumerate}
\end{problem}

\begin{solution}

\textbf{(a)} This problem has the same structure as Problem~\ref{prob:cone}, with the wire playing the role of the string. Two forces act on the bead: its weight $mg$ and the normal force $N$ exerted by the wire. Since the wire is frictionless, it can exert a force only perpendicular to itself; within the plane of the loop, this direction is along the radius, towards the centre $O$ of the loop. (The normal force could in principle also have a component perpendicular to the plane of the loop. Why is that component zero here?)

\begin{center}
\websvg{figures/tikz-1eb9bf07.svg}
\end{center}

The bead moves in a horizontal circle whose centre lies on the axis of rotation, at the height of the bead. The radius of this circle is
\begin{flalign}
r = R\sin\theta &&
\end{flalign}
The centre of the circle is therefore \emph{not} $O$. The normal force points towards $O$, but the acceleration is directed horizontally towards the axis, so only the horizontal component of $N$ provides the centripetal force.

\textbf{(b)} The normal force makes an angle $\theta$ with the vertical. (Why? Compare it with the angle at $O$.) Exactly as in \eqref{eq:vert} and \eqref{eq:horiz}, Newton's second law gives
\begin{flalign}
N\cos\theta &= mg &&\\
N\sin\theta &= m\omega^2 R\sin\theta &&
\end{flalign}
For $\sin\theta \neq 0$, it follows that $N = m\omega^2 R$ and
\begin{empheq}[box=\mathbox]{align}
\cos\theta = \frac{g}{\omega^2 R}, \qquad \omega_c = \sqrt{\frac{g}{R}}
\end{empheq}
These are the results of Problem~\ref{prob:cone}, with $L$ replaced by $R$ and $T$ by $N$. With $R = 0.10$ m,
\begin{flalign}
\omega_c = \sqrt{\frac{9.8}{0.10}} \approx 9.90 \text{ rad/s} = 9.90 \times \frac{60}{2\pi} \text{ rpm} \approx 95 \text{ rpm} &&
\end{flalign}

\textbf{(c)} Solving the result of part (b) for $\omega$ gives
\begin{flalign}
\omega = \sqrt{\frac{g}{R\cos 30^\circ}} = \sqrt{\frac{9.8}{0.10 \times 0.866}} \approx 10.6 \text{ rad/s} \approx 102 \text{ rpm} &&
\end{flalign}
An angular velocity only about 7\% above $\omega_c$ is enough to raise the bead to $30^\circ$; just above $\omega_c$, the angle increases very rapidly with $\omega$.

\textbf{(d)} No. The mass cancels, exactly as in Problem~\ref{prob:cone}. The value of 10 mg would be needed only to find the magnitude of the normal force.

\textbf{(e)$^\ast$} Consider the forces \emph{along} the wire, to which the normal force makes no contribution. Let $\hat t$ denote the unit vector along the wire in the direction of increasing $\theta$, that is, upwards and away from the axis. At angle $\theta$:
\begin{itemize}
  \item the component of gravity along $\hat t$ is $-mg\sin\theta$, which pulls the bead back towards the bottom of the loop;
  \item to remain on its horizontal circle, the bead requires an acceleration $\omega^2 R\sin\theta$ directed towards the axis, whose component along $\hat t$ is $-\omega^2 R\sin\theta\cos\theta$.
\end{itemize}
The bead can remain at angle $\theta$ only if gravity supplies exactly the required force along the wire:
\begin{flalign}
mg\sin\theta &= m\omega^2 R\sin\theta\cos\theta &&
\end{flalign}
This is the condition found in part (b), together with the root $\sin\theta = 0$.

If $\omega < \omega_c$, then $g > \omega^2 R \geq \omega^2 R\cos\theta$ for every $\theta$. Gravity therefore always pulls the bead back along the wire more strongly than its circular motion requires, and the bead slides down to the bottom from any starting position.

If $\omega > \omega_c$, then near the bottom of the loop, where $\cos\theta \approx 1$, we have $\omega^2 R\cos\theta > g$. Gravity is now too weak to hold the bead on such a small circle, so the slightest displacement causes it to move outwards until it settles at $\cos\theta = g/(\omega^2 R)$. The bottom of the loop remains a solution of the equations, but it is an unstable one, like a pencil balanced on its tip.
\end{solution}


%=====================================================================
% SOURCE: Irodov, Problems in General Physics, Sec. 1.2 (cyclist, k = k0(1 - r/R))
\begin{problem}\label{prob:cyclist}
A cyclist rides in circles on a flat horizontal platform of radius $R$. The coefficient of friction between the tyres and the platform depends only on the distance $r$ from the centre $O$ of the platform, according to
\begin{flalign}
\mu = \mu_0\br{1 - \frac{r}{R}} &&
\end{flalign}
where $\mu_0$ is a constant.

\begin{center}
\websvg{figures/tikz-c4696212.svg}
\end{center}

\begin{enumerate}[label=(\alph*)]
  \item What force keeps the cyclist moving in a circle? Is it static or kinetic friction?
  \item Find the maximum speed at which the cyclist can ride around a circle of radius $r$ centred at $O$.
  \item Find the radius of the circle on which the cyclist can ride fastest, and the corresponding maximum speed.
\end{enumerate}
\end{problem}

\begin{solution}

\textbf{(a)} The figure shows a vertical cross-section through the centre of the platform, with the forces on the cyclist (drawn as a block).

\begin{center}
\websvg{figures/tikz-6414b353.svg}
\end{center}

Gravity and the normal force are both vertical and cancel each other. The only horizontal force is the friction force $f$ exerted by the platform, which points towards $O$. This is \emph{static} friction: the part of the tyre in contact with the ground is momentarily at rest and does not slip sideways. (Why is the contact patch at rest even though the bicycle is moving?)

\textbf{(b)} Friction provides the centripetal force, and its magnitude can be at most $\mu N = \mu mg$. Therefore
\begin{flalign}
\frac{m v^2}{r} &\leq \mu_0\br{1 - \frac{r}{R}} mg \quad\Longrightarrow\quad v_{\max}^2(r) = \mu_0 g\br{r - \frac{r^2}{R}} &&
\end{flalign}

\textbf{(c)} Two effects compete here. On a larger circle, a given speed requires a smaller centripetal acceleration $v^2/r$, but the available friction decreases with distance from the centre. The function $v_{\max}^2(r)$ is a downward-opening parabola that vanishes at $r = 0$ and at $r = R$, so its maximum lies midway between these points:
\begin{empheq}[box=\mathbox]{align}
r = \frac{R}{2}, \qquad v_{\max} = \frac12\sqrt{\mu_0 g R}
\end{empheq}
The behaviour at the end points is consistent with this: at $r = 0$ the circle has no size, and at $r = R$ there is no friction at all.
\end{solution}

%=====================================================================
% SOURCE: Irodov, Problems in General Physics, Sec. 1.2 (car on the sine curve y = a sin(x/alpha))
\begin{problem}[$^\ast$]\label{prob:sine}
A car travels at constant speed along a flat horizontal road. Viewed from above, the road follows the curve $y = h\sin(x/\ell)$, where $h$ and $\ell$ are constants. The coefficient of friction between the tyres and the road is $\mu$.

\begin{center}
\websvg{figures/tikz-e50bb99f.svg}
\end{center}

\begin{enumerate}[label=(\alph*)]
  \item Given that the speed of the car is constant, in what direction does its acceleration point at each instant?
  \item At which points along the road is the car most likely to skid? Explain.
  \item Find the maximum speed at which the car can travel along the road without skidding.
\end{enumerate}
\end{problem}

\begin{solution}

\textbf{(a)} Since the speed is constant, the car has no acceleration along the road; its entire acceleration is perpendicular to the road. Over a sufficiently short stretch, any smooth curve can be approximated by an arc of a circle, whose radius $\rho$ is called the \emph{radius of curvature}. The car's acceleration then has magnitude $v^2/\rho$ and is directed towards the centre of that circle. In this sense, circular motion is a \emph{local} idea: it applies to any curved path, one short piece at a time.

Three forces act on the car: its weight $mg$ and the normal force $N$, which are vertical and cancel, and the static friction force $f$ from the road, which is horizontal and perpendicular to the road. In the top view below, only $f$ is visible; at the crest it points towards the centre of the circle that matches the road there.

\textbf{(b)} The sideways force is provided by static friction, whose magnitude is at most $\mu mg$. We therefore require $v^2/\rho \leq \mu g$ at every point of the road. This condition is most severe where the road bends most sharply, that is, where $\rho$ is smallest. For a sine curve, this occurs at the crests and troughs.

\textbf{(c)} To find $\rho$ at a crest, say at $x_0 = \pi\ell/2$, we write $x = x_0 + s$ and expand the sine function about its maximum:
\begin{flalign}
y &\approx h - \frac{h}{2\ell^2}\, s^2 &&
\end{flalign}
Near its highest point, a circle of radius $\rho$ has the form $y \approx \text{const} - s^2/(2\rho)$. (Check this.) Comparing the two expressions gives
\begin{flalign}
\rho_{\min} &= \frac{\ell^2}{h} &&
\end{flalign}
The condition $v^2/\rho_{\min} \leq \mu g$ then becomes $v^2 h/\ell^2 \leq \mu g$, so that
\begin{empheq}[box=\mathbox]{align}
v_{\max} = \ell\sqrt{\frac{\mu g}{h}}
\end{empheq}
This result is reasonable: a road with long, gentle bends (large $\ell$) or small sideways excursions (small $h$) can be driven faster.

\begin{center}
\websvg{figures/tikz-4daf2fc0.svg}
\end{center}
(In the sketch, $h = \ell$, so that $\rho_{\min} = \ell$.)
\end{solution}

\end{document}